\documentclass[11pt]{article}
\usepackage[T1]{fontenc}
\usepackage[utf8]{inputenc}
\usepackage{lmodern}
\usepackage{geometry}
\usepackage{amsmath}
\usepackage{amssymb}
\usepackage{amsthm}
\usepackage{booktabs}
\usepackage{hyperref}
\usepackage{listings}
\geometry{margin=1in}

\lstdefinestyle{lean}{
  basicstyle=\ttfamily\small,
  breaklines=true,
  columns=fullflexible,
  keepspaces=true,
  showstringspaces=false,
  frame=single,
  framesep=4pt,
  xleftmargin=2pt,
  aboveskip=8pt,
  belowskip=8pt,
  literate=
    {⟨}{{$\langle$}}1
    {⟩}{{$\rangle$}}1,
}
\lstset{style=lean}

\newtheorem{theorem}{Theorem}
\newtheorem{proposition}{Proposition}
\newtheorem{lemma}{Lemma}
\newtheorem{remark}{Remark}
\theoremstyle{definition}
\newtheorem{definition}{Definition}

\newcommand{\SDiff}{\mathrm{SDiff}}
\newcommand{\SG}{\mathrm{SG}}
\newcommand{\Ad}{\mathrm{Ad}}
\newcommand{\Hm}{H^m}
\newcommand{\Hms}{H^m_\sigma}
\newcommand{\CAd}{C_{\mathrm{Ad}}}
\newcommand{\Cphi}{C_\varphi}

\title{Machine-Checked Status of the SG--Cole--Hopf Route to\\
3D Navier--Stokes Global Regularity:\\
H1 Obstruction, Stochastic-Lagrangian Repair, and Pancake Closure}
\author{The MeTTapedia formalization effort}
\date{September 2026}

\begin{document}
\maketitle
\sloppy

\begin{abstract}
Ben Goertzel's manuscript \emph{Global Regularity of 3D Navier--Stokes via
SG--Cole--Hopf Program} (v7) reduces global regularity of the 3D incompressible
Navier--Stokes equation on $\mathbb{T}^3$ to a uniform-in-time identity-energy
bound on the volume-preserving diffeomorphism group
$\SG=\SDiff(\mathbb{T}^3)$.  The June 2026 MeTTapedia formalization found that
the route's original local adjoint hypothesis (H1) is false in a normalized
Fourier-mode shear model: inside every positive $\Hm$ chart radius there are
area-preserving shears of fixed $\Hm$ displacement size whose adjoint
$\Hm\to\Hm$ ratio is $\sqrt{1+(\varepsilon k)^2}$ and therefore unbounded.
That refutation remains intact.

This September revision records the follow-up repair campaign in full.  The
corrected route does not rely on the false $\Hm\to\Hm$ adjoint bound; it
replaces it with a stochastic-Lagrangian/cotangent transport skeleton and a
microlocal vorticity-stretching gate.  The machine-checked surface now
comprises: the evasion contrast showing the H1-breaking shear family leaves
pointwise Jacobians untouched; the one-form transport calculus with the exact
frame-Laplacian identity; the divergence-free non-exact witness that forces
the gauge; the stochastic-Lagrangian pushdown deriving Navier--Stokes with a
fully explicit pressure; the conditional BKM chain with its gate consumers;
nonzero exact inhabitants and end-to-end canaries; the frozen-strain pancake
obstruction with its exact peak-amplification law; the plane-wave
Biot--Savart null structure claimed in the pancake-closure proposal; and the
dyadic absorption step that would turn the pancake-sector estimate into a
common BKM envelope.  The September field-transfer work additionally constructs
the coherent/misaligned projection for arbitrary complex periodic Fourier
coefficients, transfers the coherent stretching estimate to those coefficients,
proves the exact square-dyadic cone lattice count, and verifies the anisotropic
$L^1$ kernel mechanism that preserves the receiver's $L^\infty$ norm.  It also
proves the estimate at continuous real frequencies and for arbitrary projected
amplitudes, constructs the exact periodic operator-kernel realization on finite
Fourier fields, transports the cancellation to every frozen oriented strain
frame, and now transports the full square-dyadic periodized kernel and its
cardinality-free $C/N$ endpoint uniformly to those frames.  Complexifying the
frame and conjugating both inputs and the output gives a physical-frame bound;
uniformity then permits an arbitrary frame choice at each output point without
a frame-derivative loss at this $L^\infty$ kernel endpoint.  A separate exact
variable/frozen split and a spectral-gap perturbation theorem remain available
for formulations that differentiate the symbol.  The physical finite-sector
stretching sum is now proved to equal this coherent annular kernel action plus
the exact three-term misalignment remainder, both in an arbitrary frozen frame
and in the full pointwise frame constructed from the strain.  This frame maps
the most expanding eigenvector to the coherent reference $e_x$ and the most
compressive eigenvector to the axial reference $e_z$.  Hence the
coherent term has a cardinality-free $C/N$ bound without an abstract split
hypothesis.  A further checked projection argument bounds the exact complex
Fourier misalignment coefficient by its energy transverse to the expanding
line.  This estimate sums over the genuine finite pancake sector without a
mode-count loss, and finite Parseval identifies the right-hand side with a
physical-space $L^2$ transverse-energy integral for a frozen full strain
frame.  A sign-invariant line-projector perturbation theorem further controls
freezing error by $4\varepsilon/\mathrm{gap}$ times the full $L^2$ energy,
with an explicit small frozen-top-gap alternative.  On the separated branch,
a homogeneous spectral-gap estimate now converts the frozen transverse
energy into the integrated top Rayleigh defect, without normalizing vorticity
at its zeros.  Finite aggregation of the coercive cells and generic
hard-set/smooth-weight bounded-overlap integral estimates are also checked.
The analytic localization bridge (including the frequency leakage produced
by spatial cutoffs), space--time treatment of the small-gap branch, and
conversion from this $L^2$/defect control to the scale-local endpoint budget
remain open.  The
latest audit also isolates why the proposed pancake closure does
not yet prove the Millennium statement: its high-low split covers all aperture
angles iff the scale-local strain satisfies $\gamma\le\nu\lambda_q$, exactly
the quantitative budget it set out to prove.  The surviving open input is the
leading route-specific pointwise-in-time Lean predicate
\texttt{ScaleLocalPancakeStrainBudget}, a sharpened version of the earlier
\texttt{StochasticStretchingEstimate}.  Field realization, the stated
non-pancake sector interfaces, and the final BKM continuation plumbing are
additional explicit obligations.  Thus the status is neither a completed
Navier--Stokes proof nor a vague ``hard analytic gap'': it is a conditional
route with a named, scale-local, machine-audited leading analytic pin and a
finite list of surrounding formalization obligations.
\end{abstract}

\section{The problem and the role of hypothesis (H1)}
\label{sec:problem}

\subsection{The SG--Cole--Hopf route}

We work on the flat three-torus $\mathbb{T}^3$ with smooth divergence-free
initial data $u_0$. Ben Goertzel's manuscript \emph{Global Regularity of 3D
Navier--Stokes via SG--Cole--Hopf Program}, version~v7 (hereafter ``v7''), lifts
the incompressible Navier--Stokes equation to a diffusion on the
volume-preserving diffeomorphism group
\[
  \SG := \SDiff(\mathbb{T}^3),
\]
equipped with a right-invariant Fourier frame $\{D_i\}_{i\ge 1}$ induced by a
divergence-free frame $\{E_i\}_{i\ge 1}$ on $\mathbb{T}^3$. The generator and
carré du champ are
\[
  L := \sum_{i\ge 1} D_i^2,
  \qquad
  \Gamma(F) := \sum_{i\ge 1} (D_i F)^2 .
\]
One solves the SG heat equation $\partial_t \Phi = \nu L\Phi$ with a strictly
positive datum $\varphi = \exp(-W_0/(2\nu))$, applies the Cole--Hopf transform
$W = -2\nu \log \Phi$ to obtain an SG Hamilton--Jacobi--Bellman potential, and
pushes the $\Gamma$-gradient down at the identity to a projected
Navier--Stokes solution $u(t)$ on $\mathbb{T}^3$.

The analytic content is concentrated entirely in closing the
Beale--Kato--Majda (BKM) vorticity criterion. With $u(t)=\mathrm{Ev}(W(t))$ and
$\omega(t)=\operatorname{curl}u(t)$, the push-down identity and Cauchy--Schwarz
give a frame constant $C_{\mathrm{curl}E}<\infty$ with
\begin{equation}
  \|\omega(t)\|_{L^\infty(\mathbb{T}^3)}
    \le \sqrt{C_{\mathrm{curl}E}}\,\sqrt{\Gamma(W(t))(\mathrm{id})},
  \label{eq:vort}
\end{equation}
so that a uniform-in-time bound on $\Gamma(W(t))(\mathrm{id})$ makes the BKM
integral $\int_0^T \|\omega(t)\|_{L^\infty}\,dt$ finite on every finite interval.
By the log chain rule $\Gamma(\log\Phi)=\Gamma(\Phi)/\Phi^2$ and the maximum
principle $\Phi(t,\mathrm{id})\ge m_\varphi$ (v7, Prop.~7.1),
\begin{equation}
  \Gamma(W(t))(\mathrm{id})
    = 4\nu^2\,\Gamma(\log\Phi(t))(\mathrm{id})
    \le \frac{4\nu^2}{m_\varphi^2}\,\Gamma(P_t\varphi)(\mathrm{id}),
  \label{eq:logbound}
\end{equation}
which reduces the whole program to a single uniform bound on
$\Gamma(P_t\varphi)(\mathrm{id})$.

\subsection{Reduction to a local energy bound, and hypothesis (H1)}

Rather than invoking a global Bakry--\'Emery $\Gamma_2$ curvature lower bound,
v7 obtains that single bound from a \emph{local} package. Let
$\kappa : U \to V \subset \Hms(\mathbb{T}^3)$ be the chart of a small
$\SDiff^m$ neighborhood of the identity ($m\ge 5$ integer, with the
$\Gamma$-Hilbert structure at $\mathrm{id}$ equivalent to the $\Hms$ norm), let
$\chi$ be the cutoff defining $W_0$ and $\varphi$, and set the cutoff region
\[
  K := \kappa^{-1}(\operatorname{supp}\chi)\subset U .
\]
Outside $K$ one has $W_0\equiv 0$, hence $\varphi\equiv 1$ and
$\Gamma(\varphi)\equiv 0$. On $K$, v7 assumes two local hypotheses (v7,
\S5.3 and \S6, displayed at Lemma~D.18):
\begin{align}
  \text{(H1) Adjoint bounded on support:} \quad
    & \exists\,\CAd<\infty : \;
      \|\Ad_g\|_{\Hm\to\Hm}\le \CAd \quad \text{for all } g\in K,
      \label{eq:H1}\\[2pt]
  \text{(H2) Finite cutoff energy:} \quad
    & \exists\,\Cphi<\infty : \;
      \Gamma(\varphi)(g)\le \Cphi^2 \quad \text{for all } g\in K.
      \label{eq:H2}
\end{align}
From \eqref{eq:H1}--\eqref{eq:H2} v7 derives the key estimate via a conjugation
differentiation identity (v7, Prop.~6.5, proved as Lemma~F.7): for all $t\ge 0$,
\begin{equation}
  \Gamma(P_t\varphi)(\mathrm{id}) \;\le\; \CAd^2\,\Cphi^2 .
  \label{eq:prop65}
\end{equation}
The mechanism is transparent. In the flow representation
$P_t F(g)=\mathbb{E}[F(G_t\circ g)]$, differentiating along a conjugation curve
$C_{t,i}(\varepsilon)=g\circ\exp(\varepsilon E_i)\circ g^{-1}$ produces the
direction $\Ad_g E_i$, and Cauchy--Schwarz in the $\Gamma$-Hilbert structure
gives the pointwise bound
\begin{equation}
  \bigl|D_{\Ad_g E_i}\varphi(g)\bigr|^2
    \le \|\Ad_g\|_{\mathrm{op}}^2\,\Gamma(\varphi)(g)
    \le \CAd^2\,\Cphi^2\,\mathbf{1}_{\{g\in K\}} ,
  \label{eq:cauchyschwarz}
\end{equation}
after which a Jensen/energy reduction yields \eqref{eq:prop65}. Chaining
\eqref{eq:prop65} through \eqref{eq:logbound} and \eqref{eq:vort} closes BKM.
Hypothesis~(H1) is therefore the single load-bearing analytic input of the
route: \emph{every} subsequent estimate inherits the constant $\CAd$ from
\eqref{eq:H1}, and the program establishes global regularity if and only if
\eqref{eq:H1} holds.

\subsection{How v7 attempts to verify (H1)}

v7 verifies~(H1) at Lemma~6.1 (and again at Lemma~D.18). The adjoint /
push-forward action is written, for $g\in\SDiff^m$ and $v\in\Hms$, as
\begin{equation}
  \Ad_g v := (Dg)\,\bigl(v\circ g^{-1}\bigr),
  \label{eq:adg}
\end{equation}
and the argument is that, for $m>5/2$, composition and inversion are continuous
(indeed smooth) operations on $\SDiff^m$, so the map $(g,v)\mapsto \Ad_g v$ is
continuous from $\SDiff^m\times\Hms$ into $\Hms$; hence
$g\mapsto\|\Ad_g\|_{\mathrm{op};\,\Hm\to\Hm}$ is locally bounded near
$\mathrm{id}$, and one shrinks $K$ until \eqref{eq:H1} holds.

The gap is in \eqref{eq:adg} itself. The factor $Dg$ is the Jacobian of $g$, and
a map $g\in\Hm$ has Jacobian $Dg\in H^{m-1}$ only. The adjoint thus carries a
\emph{derivative loss}: $\Ad_g$ maps $\Hm$ into $H^{m-1}$, and there is no
uniform $\Hm\to\Hm$ operator bound near the identity in the $\Hm$ topology.
This is the classical Ebin--Marsden phenomenon: right translation on
$\SDiff^m$ is smooth, but the adjoint of the right-invariant frame is not a
bounded $\Hm$ operator uniformly over an $\Hm$-neighborhood. The
local-boundedness argument conflates topological-group continuity of
composition with the operator-norm estimate \eqref{eq:H1} that the route
actually needs.

\section{The Fourier-mode shear model and the checked $\neg$(H1) theorems}
\label{sec:lean}

We make the derivative-loss obstruction precise and machine-checkable in a clean
analytic model, on the normalized two-torus, of an area-preserving shear family
\[
  g_k(x,y) = \bigl(x,\; y + a_k \sin(kx)\bigr),
  \qquad
  \delta_k := g_k - \mathrm{id} = \bigl(0,\; a_k\sin(kx)\bigr).
\]
The amplitude $a_k$ is chosen so that the displacement has \emph{fixed} $\Hm$
size $\varepsilon$ for every wavenumber $k$; thus the whole family sits inside
any prescribed chart radius. Applying the push-forward \eqref{eq:adg} to the
constant horizontal field $v=(1,0)$ gives
\[
  \Ad_{g_k} v = \bigl(1,\; a_k\, k\, \cos(kx)\bigr),
\]
and, with the same $\Hm$ normalization, the operator ratio is
\[
  \frac{\|\Ad_{g_k}v\|_{\Hm}}{\|v\|_{\Hm}}
    = \sqrt{1+(\varepsilon k)^2}.
\]
The displacement stays at size $\varepsilon$ while the ratio grows without bound:
no finite $\CAd$ can dominate it on any positive chart radius.

This is formalized, with no \texttt{sorry} and by elementary means, in the Lean
module \texttt{Mettapedia.FluidDynamics.NavierStokes.BenH1Break}. The model
quantities are defined as follows.

\begin{lstlisting}[language={}]
/-- In the normalized two-torus shear family, every selected displacement has
fixed `H^m` size `epsilon`. -/
def benH1ModeDisplacementHNorm (epsilon : Real) (_k : Nat) : Real :=
  epsilon

/-- Lower-bound ratio obtained by applying the shear push-forward to the
constant horizontal vector field. -/
def benH1ModeAdjointRatio (epsilon : Real) (k : Nat) : Real :=
  Real.sqrt (1 + (epsilon * (k : Real)) ^ 2)

/-- The linearized commutator part alone grows like `epsilon * k`. -/
def benH1ModeCommutatorRatio (epsilon : Real) (k : Nat) : Real :=
  epsilon * (k : Real)
\end{lstlisting}

Hypothesis~(H1), restricted to this normalized family inside a chart radius
$\rho$ (taking $\varepsilon=\rho/2$ so the family is strictly inside the ball),
is the predicate:

\begin{lstlisting}[language={}]
/-- H1 restricted to the normalized Fourier-mode shear family inside chart
radius `rho`. -/
def benH1FourierModeUniformBoundWithinRadius (rho C : Real) : Prop :=
  forall k : Nat,
    benH1ModeDisplacementHNorm (rho / 2) k <= rho ->
      benH1ModeAdjointRatio (rho / 2) k <= C
\end{lstlisting}

The core fact is that the adjoint ratio dominates the linear commutator part
$\varepsilon k$, which is unbounded over $k$. The proof is fully elementary: it
picks $k>C/\varepsilon$, lower-bounds $\sqrt{1+(\varepsilon k)^2}$ by
$\varepsilon k$, and derives a contradiction.

\begin{lstlisting}[language={}]
theorem benH1ModeAdjointRatio_unbounded {epsilon C : Real}
    (hepsilon : 0 < epsilon) :
    Not (forall k : Nat, benH1ModeAdjointRatio epsilon k <= C) := by
  intro hbound
  obtain ⟨k, hk⟩ := exists_nat_gt (C / epsilon)
  have hC_lt : C < epsilon * (k : Real) := by
    have hmul := mul_lt_mul_of_pos_left hk hepsilon
    field_simp [ne_of_gt hepsilon] at hmul
    exact hmul
  have hnonneg : 0 <= epsilon * (k : Real) := by positivity
  have hlinear_le_ratio :
      epsilon * (k : Real) <= benH1ModeAdjointRatio epsilon k := by
    unfold benH1ModeAdjointRatio
    have hsquare :
        (epsilon * (k : Real)) ^ 2 <= 1 + (epsilon * (k : Real)) ^ 2 := by
      linarith
    calc
      epsilon * (k : Real)
          = Real.sqrt ((epsilon * (k : Real)) ^ 2) := by
              rw [Real.sqrt_sq_eq_abs, abs_of_nonneg hnonneg]
      _ <= Real.sqrt (1 + (epsilon * (k : Real)) ^ 2) :=
              Real.sqrt_le_sqrt hsquare
  exact not_lt_of_ge (le_trans hlinear_le_ratio (hbound k)) hC_lt
\end{lstlisting}

From this the refutation of~(H1) on the model follows directly. Any positive
chart radius $\rho$ admits the normalized family at $\varepsilon=\rho/2$ (whose
displacements all have $\Hm$ size $\rho/2\le\rho$), and its adjoint ratios are
unbounded, contradicting any proposed finite $C$.

\begin{lstlisting}[language={}]
/-- Any positive chart radius contains the normalized shear family with
`epsilon = rho / 2`, and its adjoint ratios are unbounded. -/
theorem not_benH1FourierModeUniformBoundWithinRadius {rho C : Real}
    (hrho : 0 < rho) :
    Not (benH1FourierModeUniformBoundWithinRadius rho C) := by
  intro hbound
  have hepsilon : 0 < rho / 2 := by positivity
  have hmode_bound : forall k : Nat, benH1ModeAdjointRatio (rho / 2) k <= C := by
    intro k
    apply hbound k
    unfold benH1ModeDisplacementHNorm
    linarith
  exact benH1ModeAdjointRatio_unbounded hepsilon hmode_bound

/-- Compact public verdict: Ben's H1, as a uniform `H^m -> H^m` adjoint bound,
is false on the normalized Fourier-mode shear model for every positive chart
radius and every proposed finite constant. -/
theorem benGoertzelH1_false_in_fourierModeShearModel :
    forall rho C : Real,
      0 < rho -> Not (benH1FourierModeUniformBoundWithinRadius rho C) := by
  intro rho C hrho
  exact not_benH1FourierModeUniformBoundWithinRadius (rho := rho) (C := C) hrho
\end{lstlisting}

The three checked statements are
\texttt{benH1ModeAdjointRatio\_unbounded},
\texttt{not\_benH1FourierModeUniformBoundWithinRadius}, and the compact verdict
\texttt{benGoertzelH1\_false\_in\_fourierModeShearModel}: for every positive
chart radius and every proposed finite constant, the modeled (H1) uniform
$\Hm\to\Hm$ adjoint bound fails. (A companion lemma
\texttt{benH1ModeCommutatorRatio\_unbounded} records that the linearized
commutator part $\varepsilon k$ is already unbounded on its own.)

\section{The numeric lab: the norms are derived, not assumed}
\label{sec:lab}

The Lean development encodes the two normalized quantities
$\|\delta_k\|_{\Hm}=\varepsilon$ and
$\|\Ad_{g_k}v\|_{\Hm}/\|v\|_{\Hm}=\sqrt{1+(\varepsilon k)^2}$ as definitions. To
confirm that these are not circular --- that the factor $k$ in the adjoint ratio
is a genuine consequence of the Sobolev-weight structure and the derivative in
\eqref{eq:adg}, rather than an assumption --- a companion numeric lab
(\texttt{ns\_ben\_h1\_fourier\_mode\_lab.py}, in the MeTTapedia validation lab)
\emph{computes} the norms from first quantities. With $\Hm$ Fourier weight
$w_k = (1+k^2)^{m/2}$:
\begin{itemize}
  \item the amplitude is set to $a_k = \sqrt{2}\,\varepsilon / w_k$, so the
        displacement norm is
        $\|\delta_k\|_{\Hm} = |a_k|\,w_k/\sqrt{2} = \varepsilon$
        (the $\sqrt2$ is the $L^2$ normalization of $\sin(kx)$);
  \item differentiating, $\partial_x\bigl(a_k\sin(kx)\bigr)=a_k k\cos(kx)$, whose
        $\Hm$ norm is $|a_k|\,k\,w_k/\sqrt{2} = \varepsilon k$. The weight $w_k$
        \emph{cancels} the amplitude, leaving $\varepsilon$; the surviving factor
        $k$ is the one derivative gained by $D g_k$;
  \item the constant field has $\|v\|_{\Hm}=1$, so the adjoint norm is
        $\sqrt{1^2+(\varepsilon k)^2}$ and the ratio is
        $\sqrt{1+(\varepsilon k)^2}$.
\end{itemize}
This is exactly the Ebin--Marsden derivative loss of \eqref{eq:adg} made
quantitative: $g_k\in\Hm$ has $Dg_k\in H^{m-1}$, and the action of $Dg_k$
multiplies the mode's high-frequency content by its wavenumber. Table~\ref{tab:lab}
lists the lab output for $\varepsilon=0.01$, $m=3$. The displacement norm stays
fixed at $\varepsilon$ across four decades of $k$ while the adjoint ratio grows.

\begin{table}[h]
\centering
\begin{tabular}{rcccc}
\toprule
$k$ & amplitude $a_k$ & $\|\delta_k\|_{\Hm}$ & commutator $\varepsilon k$ & ratio $\sqrt{1+(\varepsilon k)^2}$ \\
\midrule
$8$    & $2.699\times10^{-5}$ & $0.0100$ & $0.080$ & $1.0032$ \\
$16$   & $3.433\times10^{-6}$ & $0.0100$ & $0.160$ & $1.0127$ \\
$32$   & $4.310\times10^{-7}$ & $0.0100$ & $0.320$ & $1.0500$ \\
$64$   & $5.393\times10^{-8}$ & $0.0100$ & $0.640$ & $1.1873$ \\
$128$  & $6.743\times10^{-9}$ & $0.0100$ & $1.280$ & $1.6243$ \\
$256$  & $8.429\times10^{-10}$ & $0.0100$ & $2.560$ & $2.7484$ \\
$512$  & $1.054\times10^{-10}$ & $0.0100$ & $5.120$ & $5.2167$ \\
$1024$ & $1.317\times10^{-11}$ & $0.0100$ & $10.240$ & $10.2887$ \\
\bottomrule
\end{tabular}
\caption{Fourier-mode lab output ($\varepsilon=0.01$, $m=3$). The displacement
$\Hm$ norm is pinned at $\varepsilon$ for every wavenumber, while the adjoint
operator ratio $\sqrt{1+(\varepsilon k)^2}$ grows without bound. For any
proposed constant $C$ the ratio first exceeds it near
$k\approx\sqrt{C^2-1}/\varepsilon$ (e.g.\ $C=2$ at $k=174$, $C=10$ at $k=995$,
$C=100$ at $k=10000$): no finite $\CAd$ survives.}
\label{tab:lab}
\end{table}

\section{The repaired stochastic-Lagrangian route: the full formalized layer}
\label{sec:repair}

After the H1 obstruction was isolated, the formalization effort tested a
different repair strategy.  The repaired route abandons the false
$\Hm\to\Hm$ adjoint estimate and moves the load-bearing estimate to the
stochastic-Lagrangian derivative flow of Constantin--Iyer (\emph{Comm.\ Pure
Appl.\ Math.}\ \textbf{61} (2008), 330--345): the state is a momentum
one-form $\mu = u^\flat + dW$ transported along a noisy flow, the velocity is
recovered by averaging and projection, and the continuation estimate is a
vorticity-stretching gate feeding the BKM criterion.  In the Lean development
the gate appears first as the coarse input
\texttt{StochasticStretchingEstimate}, and in the dyadic refinement as the
scale-local predicate \texttt{ScaleLocalPancakeStrainBudget}.

This section records the \emph{whole} machine-checked surface of the repair
campaign, layer by layer, not only the current crux.  All Lean listings below
are structurally verbatim from the sources (declaration names, hypotheses,
and conclusions exactly as checked), with Unicode notation transliterated to
ASCII for typesetting; the Lean sources are normative.  Table~\ref{tab:surface}
summarizes the layers.

\begin{table}[h]
\centering
\small
\begin{tabular}{lll}
\toprule
Layer (module) & Headline results & Scope label \\
\midrule
H1 model (\texttt{BenH1Break}) & adjoint ratio unbounded; $\neg$(H1) & model refutation \\
Evasion (\texttt{ShearJacobianBounds}) & same family: pointwise Jacobian $\le\varepsilon$ & model contrast \\
One-forms (\texttt{OneFormFrameCalculus}) & frame-Laplacian identity $\sum_i L^2_{E_i}=\Delta$ & exact identity \\
Exactness (\texttt{OneFormExactness}) & div-free smooth non-exact witness & exact witness \\
Gradients (\texttt{GradientOneForm}) & Cartan formula; self-transport identity & exact identities \\
Pushdown (\texttt{ConstantinIyerRepresentation}) & NS with explicit pressure from transport & conditional (interface) \\
Instances (\texttt{DecayingShearInstance}, regression) & nonzero exact inhabitant; rest canaries & exact instances \\
Gate (\texttt{VorticityStretchingGate}) & pointwise common envelope from named pin & conditional \\
Countermodel (\texttt{FrozenStrainModel}) & exact pancake solution; no uniform depletion & model refutation \\
Aperture (\texttt{ApertureCoverageDichotomy}) & dichotomy $\iff$ strain budget; gap failures & model equivalence \\
Null structure (\texttt{PlaneWaveSelfStrain}) & single-mode Biot--Savart null stretching & exact identity \\
Dyadic closure (\texttt{DyadicPancakeClosure}) & absorption algebra; sharpened pin; route & conditional \\
Coherence (\texttt{PancakeCoherentPairEstimate}) & pairwise aperture gain & exact inequality \\
Refined pin (\texttt{MisalignmentRefinedPin}) & misalignment budget implies pancake budget & conditional \\
Direction dynamics (\texttt{DirectionEvolutionTilting}) & direction laws; gap-free tilting no-go & exact/refutation \\
\bottomrule
\end{tabular}
\caption{The formalized surface of the repair campaign (all modules under
\texttt{Mettapedia.FluidDynamics.NavierStokes}, September 2026).  ``Model''
statements are proved in named analytic models; ``exact'' statements are
unconditional identities or witnesses on $\mathbb{R}^3$-coordinate space;
``conditional'' statements are machine-checked implications whose analytic
inputs are named hypotheses.}
\label{tab:surface}
\end{table}

\subsection{Why a pointwise transport route can evade the H1 countermodel}
\label{sec:evasion}

The first question for any repair is whether the Section~\ref{sec:lean}
countermodel already defeats it.  It does not, and this is itself a theorem.
On \emph{the same} $\Hm$-normalized shear family that refutes (H1), the
module \texttt{ShearJacobianBounds} proves
(\texttt{shearFamily\_pointwise\_bounded\_while\_sobolev\_ratio\_unbounded}):

\begin{enumerate}
  \item every pointwise Jacobian deviation is bounded by $\varepsilon$
        uniformly in the frequency $k$
        (\texttt{shearFamily\_jacobian\_uniform\_bound}:
        $\|Dg_k(x)v - v\| \le \varepsilon\|v\|$ for all $x,v$); while
  \item the $\Hm\to\Hm$ operator ratio $\sqrt{1+(\varepsilon k)^2}$ of the
        same family admits no finite uniform bound (Section~\ref{sec:lean}).
\end{enumerate}

The high-frequency-shear derivative-loss mechanism therefore lives entirely
in the Sobolev operator topology: the quantity feeding a pointwise gate
$\int_0^T \mathbb{E}\|DX_t\|_\infty\,dt$ is not stressed at all by this
family.  The checked statement is deliberately labeled: it does not prove any
pointwise gate --- it proves this family cannot refute one.  That is what
licenses moving the route from adjoint bounds to Lagrangian transport.

\subsection{The one-form layer: non-exact momenta and the frame Laplacian}
\label{sec:oneform}

The repaired state is a momentum one-form with a gauge, $\mu=u^\flat+dW$,
not a scalar potential.  Two checked facts pin down why this shape is forced
and why the frame generator acts correctly on it.

\emph{Non-exact momenta are necessary.}  Exact one-forms are closed
(\texttt{oneFormExtDerivAt\_eq\_zero\_of\_exact}), and there is a smooth,
everywhere divergence-free field that is not exact --- the transverse sine
mode:

\begin{lstlisting}[language={}]
theorem exists_divergenceFree_smooth_not_exact_dim3 :
    exists u0 : EuclideanSpace Real (Fin 3) -> EuclideanSpace Real (Fin 3),
      ContDiff Real 2 u0 /\
      (forall x, coordDivergence u0 x = 0) /\
      Not (IsExactOneForm u0)
\end{lstlisting}

A scalar-potential (pure Cole--Hopf) ansatz therefore cannot carry a general
divergence-free momentum: vorticity lives exactly in the non-exact part.
This is the constructive face of the Section~\ref{sec:problem} diagnosis, and
it is what the additive gauge $dW$ is for.

\emph{The frame reproduces the Laplacian.}  The route's generator is a sum of
squared frame derivatives (v7's $L=\sum_i D_i^2$).  On the truncated
trigonometric coordinate frame (three axes, transverse polarizations, both
phases) the squared Lie derivatives of a one-form sum \emph{exactly} to the
coordinate Laplacian: with amplitude $1/\sqrt2$ in dimension~3,
\[
  \sum_{j}\sum_{p\ne j}
    \Bigl(L^2_{\cos\text{-mode}(j,p)} + L^2_{\sin\text{-mode}(j,p)}\Bigr)\mu
  = \Delta\mu
\]
(\texttt{frameLaplacian\_oneForm\_normalized}, from the general
$(\mathrm{card}\,\iota-1)a^2$ identity \texttt{frameLaplacian\_oneForm} and
the phase-pair cancellation \texttt{lieOneForm\_sq\_cos\_add\_sin}).  The
supporting calculus (\texttt{GradientOneForm}) includes the coordinate Cartan
formula for exact forms (\texttt{lieOneForm\_gradField}: the Lie derivative
of $dV$ along $u$ is exact with potential $\langle u,\nabla V\rangle$, i.e.\
$L_u\,dV = d\langle u,\nabla V\rangle$ pointwise) and the self-transport
identity (\texttt{lieOneForm\_self}: $L_u u^\flat = (u\cdot\nabla)u +
\tfrac12\nabla|u|^2$).  These two identities are what turn transport
equations into Navier--Stokes with a pressure, next.

\subsection{The pushdown theorem: Navier--Stokes with an explicit pressure}
\label{sec:pushdown}

The structure \texttt{TransportedMomentumData} packages a velocity $u$, a
gauge $W$, their certified time derivatives, the gradient/Laplacian exchange
facts, and one transport hypothesis --- the It\^o-averaged frame-transport
equation for the momentum one-form:
\[
  \partial_t u + \nabla \dot W
    + L_u\bigl(u^\flat + dW\bigr)
  = \nu \,\Sigma_{\mathrm{frame}}\bigl(u^\flat + dW\bigr),
\]
where $\Sigma_{\mathrm{frame}}$ is the normalized frame second-order sum of
Section~\ref{sec:oneform} (the hypothesis is stated in averaged form; no
stochastic calculus is formalized, and the field is labeled accordingly in
the source).  From this single hypothesis the pushdown theorem derives the
Navier--Stokes equations with a fully explicit pressure:

\begin{lstlisting}[language={}]
/-- Pushdown theorem: the averaged frame-transport equation for the
momentum one-form u^flat + dW forces u to satisfy the Navier-Stokes
equations with the explicit pressure
p = dW/dt + (1/2)<u,u> + <u, grad W> - nu * Lap W. -/
theorem navierStokes_of_transportedMomentum (t : Real) (x : R3) :
    deriv (fun s => D.u s x) t + fderiv Real (D.u t) x (D.u t x) +
      gradField (D.pressure t) x =
    D.nu smul coordLaplacian (D.u t) x
\end{lstlisting}

with
\[
  p \;=\; \partial_t W \;+\; \tfrac12\langle u,u\rangle
      \;+\; \langle u,\nabla W\rangle \;-\; \nu\,\Delta W .
\]
The pressure is not posited: it is assembled from the exact terms of the
transport (gauge drift, kinetic energy via the self-transport identity, gauge
advection via the Cartan formula, viscous gauge term), and the theorem is
unconditional given the data.  This is the repaired route's replacement for
v7's identity-energy pushdown: the same endpoint (a genuine Navier--Stokes
solution with pressure), reached without the false adjoint bound.

\subsection{The vorticity-stretching gate and the conditional route}
\label{sec:gate}

Continuation is fed by the averaged Cauchy formula
$\omega(t)=\mathbb{E}[(DX_t)(\omega_0\circ A_t)]$, which bounds
$\|\omega(t)\|_\infty$ by moments of the flow Jacobian.  The gate is the
named open input of the coarse route, exactly:

\begin{lstlisting}[language={}]
/-- The vorticity-stretching gate (named analytic input; OPEN). -/
def StochasticStretchingEstimate
    (D : StochasticCauchyVorticityData Omega) (T : Real) : Prop :=
  IntervalIntegrable D.expectedJacobianBound volume 0 T
\end{lstlisting}

Its consumers are mechanical and checked: a time-uniform Jacobian bound gives
a time-uniform vorticity bound
(\texttt{vorticity\_uniform\_bound\_of\_jacobian\_bound}), and under the gate
the BKM integrand $t\mapsto\|\omega(t,x)\|$ is interval-integrable with
$\int_0^T\|\omega\| \le \|\omega_0\|_\infty \int_0^T
\mathbb{E}\|DX_t\|_\infty$
(\texttt{bkm\_integrand\_intervalIntegrable}, \texttt{bkm\_integral\_bound}).
The assembled conditional route
(\texttt{stochasticLagrangian\_conditional\_route}) states: given transported
momentum data, averaged-Cauchy vorticity data linked by
$\omega=\operatorname{curl}u$, and the gate on $[0,T]$, the velocity
satisfies Navier--Stokes with the explicit pressure \emph{and} the BKM
integrand is controlled.  Every step is machine-checked; the only analytic
input is the gate.

\emph{The route is inhabited.}  Non-vacuity is itself checked, twice.  The
rest state instantiates the route end-to-end
(\texttt{conditional\_route\_rest}, with the gate discharged by a constant
envelope).  Beyond the rest state, the module \texttt{DecayingShearInstance}
provides a nonzero exact solution of the full transported-momentum system:
the decaying shear $u(t,y)=e^{-t}\sin(y_1)\,e_0$ with gauge amplitude
$b(t)=(e^{-2t}-e^{-4t})/8$ on the frequency-doubled mode, whose balance
$b'=-4b+\tfrac14 e^{-2t}$ (\texttt{gaugeAmp\_balance}) absorbs the kinetic
self-transport term; \texttt{decayingShear\_navierStokes} confirms the pair
satisfies Navier--Stokes with the representation pressure at $\nu=1$.  So the
conditional machinery is not vacuously true: its hypotheses have nonzero
exact witnesses.

\subsection{The frozen-strain countermodel and the aperture geometry}
\label{sec:frozen}

Any proof of the gate must use more than local strain geometry plus
stochastic averaging.  The frozen-strain model makes this a theorem.  For the
prescribed strain $A=\mathrm{diag}(2\gamma,-\gamma,-\gamma)$ and the pancake
packet (vorticity along the expanding axis, wave vector along a compressive
axis), the module \texttt{FrozenStrainModel} proves that
\[
  \omega(t,x) = a(t)\,\sin\bigl(k(t)\,x_1\bigr)\,e_0,
  \qquad
  k(t)=k_0e^{\gamma t},
  \qquad
  a(t)=a_0\exp\Bigl(2\gamma t - \nu k_0^2\tfrac{e^{2\gamma t}-1}{2\gamma}\Bigr)
\]
is an \emph{exact} solution of the frozen vorticity equation
$\partial_t\omega+(Ax\cdot\nabla)\omega=A\omega+\nu\Delta\omega$
(\texttt{frozenStrainVorticity\_solves}; the drift, stretching, phase and
viscous terms are all checked, not modeled).  Writing
$g=2\gamma/(\nu k_0^2)$ for the normalized strain, the amplification peaks at
$t^*=\log g/(2\gamma)$ with the exact value
\[
  \frac{a(t^*)}{a_0} = g\,e^{1/g-1}
  \qquad(\texttt{frozenStrainAmplification\_peak}),
\]
which is unbounded over the strain family: no uniform depletion bound of the
form ``pancake alignment always acquires diffusive decay at rate
$\nu\lambda^2$'' survives
(\texttt{frozenStrain\_no\_uniform\_amplification\_bound}).  Any correct gate
proof must therefore use the \emph{self-consistency} of the Navier--Stokes
strain --- a statement quantified over arbitrary imposed strains is false.

The angular version is the aperture-coverage geometry.  Anisotropic
Bernstein estimates control cone apertures up to $\lambda^{-1/2}$;
stretch--diffusion covariance controls apertures above
$\sqrt{\gamma/\nu}/\lambda$.  The band claimed by neither,
\[
  \bigl(\lambda^{-1/2},\; \sqrt{\gamma/\nu}/\lambda\bigr),
  \qquad\text{is nonempty}\iff \nu\lambda<\gamma
  \qquad(\texttt{anisotropic\_coverage\_gap}).
\]

\subsection{July additions: the aperture dichotomy is the strain budget}
\label{sec:aperture}

The July campaign sharpened the coverage gap into an equivalence and closed
the remaining escape routes.  In \texttt{ApertureCoverageDichotomy.lean}:

\begin{lstlisting}[language={}]
theorem pancakeApertureCovered_iff_strain_le
    (hnu : 0 < nu) (hlam : 0 < lam) :
    (forall theta : Real,
        theta <= (Real.sqrt lam)^{-1} \/ Real.sqrt (gamma / nu) / lam <= theta)
      <-> gamma <= nu * lam
\end{lstlisting}

That is: the ``narrow $\Rightarrow$ enstrophy-controlled, else
$\Rightarrow$ angularly coercive'' dichotomy tiles the whole aperture range
\emph{exactly when} the scale-local strain budget $\gamma\le\nu\lambda_q$
holds.  On the gap band both mechanisms fail quantitatively:
\texttt{gapAperture\_diffusionRate\_lt\_strain} proves
$\nu(\theta\lambda)^2<\gamma$ below the stretch--diffusion threshold (the
diffusive rate is strictly beaten by the strain), and
\texttt{gapAperture\_bernstein\_volume\_exceeds\_planar\_budget} proves
$\lambda^2<\theta^2\lambda^3$ above the Bernstein threshold (the cone's
Fourier-support volume exceeds the planar budget behind the improved
$L^2\to L^\infty$ constant).

The time-averaged escape route is closed by a Gr\"onwall-saturation
computation: over the coherent window the packet consumes an integrated
strain budget of exactly $\int_0^{t^*}2\gamma\,dt=\log g$
(\texttt{frozenStrain\_integratedStrain\_peakWindow}) while amplifying to
$g\,e^{1/g-1}\ge e^{\log g-1}$, i.e.\ amplification at least
$e^{\,\mathrm{budget}-1}$
(\texttt{frozenStrainAmplification\_peak\_ge\_exp\_integratedStrain\_sub\_one}).
An averaged strain hypothesis therefore bounds amplification only
exponentially in itself and cannot substitute for the pointwise-in-time
aperture condition in a dyadic relaxation estimate.

\subsection{The dyadic pancake closure: verified pieces and the isolated pin}
\label{sec:pancake}

The July repair proposal claims a complete closure of the gate: decompose the
weighted stretching sum by angular sector, prove the pancake-sector estimate
\[
  \int_0^T \sum_{q\ge Q}\lambda_q^{-2}
    \bigl\|A_q\,\omega_q^{\mathrm{pan}}\bigr\|_\infty\,dt
  \;\le\; C(T,\nu,u_0,Q) + \varepsilon\int_0^T D_Q(t)\,dt,
  \qquad
  D_Q(t):=\sum_{q\ge Q}\|\Delta_q\omega(t)\|_\infty,
\]
with $A_q=S(S_{q-L}u)$ the low-frequency strain and
$\omega_q^{\mathrm{pan}}$ the projection onto the pancake cone of aperture
$\delta_q=\lambda_q^{-1/2}$, then absorb $\varepsilon\int D_Q$ and conclude
$\int_0^T D_Q<\infty$, hence BKM.  The formalization adjudicates this claim
piece by piece.

\emph{The genuinely new ingredient is verified.}  The proposal's Step~1 --- a
pancake vorticity mode does not generate the strain that stretches itself ---
is checked exactly, not at symbol level.  In \texttt{PlaneWaveSelfStrain}:
for $\omega(x)=\alpha\sin(\kappa x_1)e_0$ with its Biot--Savart velocity (the
shear $u(x)=-(\alpha/\kappa)\cos(\kappa x_1)e_2$),
\texttt{curlField\_planeWaveVelocityMode} proves
$\operatorname{curl}u=\omega$ exactly, and
\texttt{strainRateApply\_planeWaveVelocityMode\_vorticity} proves
$S(u)\,\omega = \tfrac12\bigl(\nabla u+(\nabla u)^{\mathsf T}\bigr)\omega=0$
pointwise, with a non-vacuity witness
(\texttt{planeWaveVorticityMode\_ne\_zero}).  The scope is honest: the
identity is for one exact mode; a pancake \emph{sector} aggregates a cone of
modes, and charging the cross-mode strain is precisely the open budget below.

\emph{The absorption skeleton is checked with explicit constants.}  The
structure \texttt{DyadicVorticityStretchingData} carries the four time
functionals ($D_Q$; pancake, angularly coercive, and residual stretching
sums), their nonnegativity and integrability, and the transported dyadic
relaxation inequality.  The three sector estimates are explicit predicates,
e.g.
\begin{lstlisting}[language={}]
def PancakeSectorEstimate (C eps : Real) : Prop :=
  (Integral of D.pancakeStretchSum over 0..D.T) <=
    C + eps * (Integral of D.highBlockSum over 0..D.T)
\end{lstlisting}
(coercive and residual analogous; both are \emph{interfaces} here ---
established in the source analysis by good-unknown/stretch--diffusion and
Coifman--Meyer arguments respectively, consumed as named hypotheses, not
re-proved).  The absorption theorem
\texttt{highBlockIntegral\_le\_of\_sectorEstimates} proves: sector estimates
with total gain $\varepsilon_p+\varepsilon_c+\varepsilon_r<1$ imply
\[
  \int_0^T D_Q(t)\,dt \;\le\;
  \frac{C_{\mathrm{relax}} + C_p + C_c + C_r}
       {1-(\varepsilon_p+\varepsilon_c+\varepsilon_r)} ,
\]
with the a-priori finiteness carried by the integrability fields (supplied in
the analytic setting by local-in-time regularity and continuation, which is
standard and not re-derived).  The theorem
\texttt{pancakeSectorEstimate\_of\_scaleLocalStrainBudget} derives the
pancake-sector estimate from the pointwise scale-local budget by time
integration.

\emph{Where the claimed proof fails.}  The proposal's high-low split charges
the whole non-pancake complement to the angularly coercive estimate.  By
Section~\ref{sec:aperture} that charge is \emph{equivalent} to the
scale-local budget $\gamma\le\nu\lambda_q$ --- which is exactly the bound the
proposal's own conditional section had identified as ``the remaining
burden'' before the claimed proof.  The claimed proof therefore consumes, at
its high-low split, the statement it set out to prove; the closure is
circular at precisely that point, and nothing else in it repairs the circle.

The remaining pin is a precise predicate, structurally verbatim:
\begin{lstlisting}[language={}]
/-- The scale-local pancake strain budget (named analytic input; OPEN --
the first load-bearing missing lemma of the claimed pancake closure). -/
def ScaleLocalPancakeStrainBudget (D : DyadicVorticityStretchingData)
    (B eps : Real) : Prop :=
  exists budget : Real -> Real,
    IntervalIntegrable budget volume 0 D.T /\
    (Integral of budget over 0..D.T) <= B /\
    forall t in Set.Icc 0 D.T,
      D.pancakeStretchSum t <= budget t + eps * D.highBlockSum t
\end{lstlisting}
Here \texttt{pancakeStretchSum} represents
$\sum_{q\ge Q}\lambda_q^{-2}\|A_q\omega_q^{\mathrm{pan}}(t)\|_\infty$.  The
obligation is to establish this pointwise envelope for the
\emph{self-consistent} Navier--Stokes strain (generated from the vorticity by
the Biot--Savart law), not for an externally imposed frozen strain --- via a
genuinely Navier--Stokes-specific mechanism (pressure/Hessian feedback,
multi-scale Biot--Savart self-depletion, or frequency-cascade accounting, in
the proposal's own taxonomy).  The verified plane-wave null structure is
evidence in the right direction; the aperture equivalence states exactly what
quantitative multi-scale estimate is still missing.

\emph{The sharpened conditional route.}  The headline theorem
\texttt{stochasticLagrangian\_conditional\_route\_of\_pancakeBudget} closes
the loop: given transported momentum data, dyadic relaxation data whose block
sum dominates the high-frequency part of $\|\operatorname{curl}u\|$
(the Littlewood--Paley decomposition interface), the scale-local pancake
strain budget, and the coercive/residual interfaces with total gain below
$1$, the velocity satisfies Navier--Stokes with the explicit pressure and,
for each fixed spatial point, the vorticity norm is interval-integrable with
an $x$-independent bound.  This reaches the same pointwise endpoint as the
coarse conditional route with the gate replaced by the sharper geometric
pin; the two pins are cross-referenced in the sources and both remain open.
The common time majorant is the correct precursor to the BKM criterion, but
the current theorem does not yet define the spatial essential supremum,
prove its measurability and integrability, or invoke a formal continuation
theorem.  An end-to-end rest-state canary
(\texttt{rest\_conditional\_route}) checks that the refined hypotheses are
simultaneously satisfiable.

\section{The direction-coherence refinement of the pin}
\label{sec:coherence}

The September 2026 campaign adjudicated the Biot--Savart self-depletion
mechanism of the pancake budget by pre-registered kinematic experiment,
then distilled and proved the surviving content, and committed a
strictly sharper pin.  All statements below are machine-checked with
the standard three axioms; the labs are
\texttt{ns\_pancake\_kinematic\_depletion\_lab.py} and
\texttt{ns\_pancake\_coherent\_aggregate\_lab.py} in the validation lab,
with pre-registered verdict criteria recorded in the files before
execution.

\emph{The dichotomy.}  Wavevector-cone confinement alone does not
deplete pairwise stretching.  In the exact lattice interaction model
(the amplitude-level Fourier symbol of $S(u)\omega$ with Biot--Savart
velocity), the aggregate cone gain equals the cone measure (measured
$D \sim \delta^{2.06}$ across dyadic shells), while the worst pair
saturates the generic maximum at every aperture.  The distilled exact
statement (\texttt{PancakeConeDirectionDichotomy}): the axial source
$k_1=(0,0,m)$ with vorticity $e_x$ stretches the \emph{perpendicular}
in-plane vorticity direction $e_y$ of the cone receiver $k_2=(1,0,M)$
with $\sigma^2 = 1/4$ exactly, for all $m, M$; the \emph{aligned}
polarization $(-M,0,1)$ of the same wavevector pair is depleted at
exactly $\sigma^2 = 1/(4(1+M^2))$.  The operative small parameter of
the self-depletion mechanism is therefore the vorticity-direction
misalignment inside the pancake --- the Constantin--Fefferman coherence
quantity --- not the wavevector cone.

\emph{The coherent pair estimate (proved).}  Give every mode the
direction-coherent polarization: the fixed in-plane reference $e_x$
projected off the wavevector.  Then
(\texttt{PancakeCoherentPairEstimate.sigmaSq\_coherent\_le}) for
wavevectors inside pancake cones of squared apertures
$d_1, d_2 \le 1/4$,
\[
  \sigma^2 \le 6\,(d_1 + d_2),
\]
with the mechanism carried by the exact identity
$k \times (|k|^2 e_x - (k\cdot e_x)k) = |k|^2 (0, k_z, -k_y)$: the
Biot--Savart velocity of a coherent mode is perpendicular to the
reference direction.  The bound aggregates over any finite family of
coherent sources (\texttt{sum\_sigmaSq\_coherent\_le}), matching the
measured coherent-aggregate gain of the second lab.

\emph{Amplitude and field-transfer layer (proved, but not yet complete).}
The module \texttt{PancakeCoherentWeightedTransfer} inserts the actual scalar
Fourier amplitudes into the pair estimate.  It proves the cardinality-free
finite convolution bound
\[
  \Bigl\|\sum_{k\in E} T_k(a_kp_k,bp_\ell)\Bigr\|_{\ell^1(\mathbb C^3)}
  \le 6(\delta_1+\delta_2)\Bigl(\sum_{k\in E}|a_k|\Bigr)|b|,
\]
and its Cauchy--Schwarz energy form, where $p_k=P_ke_0$.  The exact theorem
\texttt{PancakeConeLatticeCount.card\_le\_pancakeBox} shows that, at
$\lambda=N^2$ and aperture $1/N$, the cone lies in a box of cardinality
$(4N+1)^2(4N^2+1)=O(\lambda^2)$, rather than the ambient
$O(\lambda^3)$ count.

The modules \texttt{PancakePeriodicCoherentSplit} and
\texttt{PancakePeriodicComplexStretch} then move from the rational symbol
model to the genuine Fourier lattice $\mathbb Z^3$ and arbitrary complex
coefficients.  For every cone mode they construct the exact orthogonal split
\[
 \widehat\omega(k)=
 \frac{\langle P_ke_0,\widehat\omega(k)\rangle}
      {\langle P_ke_0,P_ke_0\rangle}P_ke_0
 +\widehat\omega_{\rm mis}(k),
 \qquad
 \langle P_ke_0,\widehat\omega_{\rm mis}(k)\rangle=0,
\]
prove that both terms remain divergence-free, define the complexified
Biot--Savart strain symbol, and prove its agreement with the rational symbol.
The aperture estimate and the $O(\lambda^2)$ energy bound therefore hold for
arbitrary complex coherent amplitudes, not only rational test vectors.

The module \texttt{PancakeRealCoherentPairEstimate} removes the rational-grid
restriction needed for a smooth multiplier.  It proves the same cancellation
for all real frequencies, constructs the rank-one orthogonal coherent
projection of an arbitrary amplitude, proves that this projection does not
increase Euclidean energy, and obtains the field-facing pointwise estimate
against the energies of the original input amplitudes.

The first attempted receiver-only endpoint has now been fenced off rather
than promoted.  The generic modules \texttt{PancakeConeKernelScaling} and
\texttt{PancakeConeOperatorKernel} correctly prove that
$K_\delta(x_\perp,z)=\delta^3K(\delta x_\perp,z)$ has
$\|K_\delta\|_{L^1}=\delta\|K\|_{L^1}$, but one of those three powers must
already be a genuine symbol gain; only two come from the transverse
Jacobian.  The periodic theorem
\texttt{PancakePeriodicOperatorKernel} now proves the corresponding no-go:
if an exact projection fixes a unit Fourier mode, its kernel has $L^1$ mass
at least one.  Thus the smooth family in
\texttt{PancakeSmoothCoherentAperture}, whose symbol was explicitly
attenuated by an additional factor $\delta$, is a useful analytic model but
is not the desired coherent field projection.

The genuine aperture factor has instead been located in the complete
bilinear interaction.  For
$k_1=(\delta a,\delta b,c)$ and
$k_2=(\delta d,\delta e,f)$,
\texttt{PancakeCoherentStretchFactorization} proves exactly, both for the
canonical coherent vectors and for coherent projections of arbitrary input
amplitudes,
\[
 T(k_1,k_2)=\delta\,T_{\rm norm}(\delta,a,b,c,d,e,f),
\]
where $T_{\rm norm}$ is smooth through $\delta=0$ before the projection
denominators are introduced.  The module
\texttt{PancakeSmoothCoherentStretchSymbol} supplies nested cutoffs keeping
both axial frequencies away from zero, constructs all $27$ entries of the
localized normalized tensor as compactly supported Schwartz symbols on six
frequency variables, proves the tensor is nonzero, and gives an integrable
inverse-Fourier kernel with exact entrywise and tensor Fourier recovery for
each fixed $\delta$.  Finally,
\texttt{PancakeBilinearKernelScaling} proves the correct four-transverse
variable scaling
\[
 K_\delta(y_\perp,z_\perp;y_\parallel,z_\parallel)
 =\delta^5K(\delta y_\perp,\delta z_\perp;
             y_\parallel,z_\parallel),
\]
whose four-dimensional Jacobian leaves one power of $\delta$, together with
the resulting bilinear $L^\infty\times L^\infty\to L^\infty$ Young bound.
The module \texttt{PancakeUniformKernelFamily} differentiates the symbol
jointly in $\delta$ and all six frequency variables, uses its fixed compact
support to obtain uniform Schwartz seminorms, converts those bounds into
uniform polynomial decay of every inverse-Fourier entry, and sums the $27$
entries to prove
\[
 \sup_{0\leq\delta\leq 1/2}\int_{\mathbb R^6}
   \lVert K_{\mathrm{norm},\delta}(y,z)\rVert\,dy\,dz < \infty.
\]
Thus the normalized Euclidean kernel-mass obligation is discharged.  The
module \texttt{PancakeBilinearPeriodization} then carries this kernel to the
standard six-torus.  It proves the fundamental-domain unfolding lemma,
matches the torus characters to the Euclidean Fourier normalization, and
shows for every $n\in\mathbb Z^6$ that
\[
 \widehat{K^{\rm per}_{\mathrm{norm},\delta}}(n)
   =M_{\mathrm{norm},\delta}(n).
\]
It also proves that the resulting six-dimensional kernel acts on a pair of
three-torus Fourier monomials through this exact coefficient.  Restoring the
algebraically extracted factor gives a periodic kernel realizing the
localized coherent-stretching symbol and one constant $C$, independent of
$0\leq\delta\leq 1/2$, such that
\[
 \|B_\delta(f,g)\|_{L^\infty(\mathbb T^3)}
 \leq 9\,\delta C\,
   \|f\|_{L^\infty(\mathbb T^3)}
   \|g\|_{L^\infty(\mathbb T^3)}.
\]
Thus Euclidean-to-torus transference and the frequency-uniform frozen-axis
coherent aperture bound are discharged.  The module
\texttt{PancakeFiniteFourierBilinearTransfer} distributes the bilinear action
over arbitrary finite trigonometric polynomials and proves the exact double
Fourier sum without a mode-cardinality loss.  It also identifies the assembled
27-entry symbol, on the retained chart, with the complex-bilinear extension
of the coherently projected normalized Biot--Savart stretching interaction;
off that chart both sides vanish.  Thus the periodic kernel now acts exactly
as the localized coherent PDE interaction on finite Fourier fields.

The module \texttt{PancakeDyadicKernelScaling} carries this base chart to
actual square-dyadic shells.  For every positive $r$ it proves
\[
 \int_{\mathbb R^6}\!\|r^6K(rx)\|\,dx
   =\int_{\mathbb R^6}\!\|K(x)\|\,dx,
 \qquad
 \widehat{r^6K(r\,\cdot)}(\xi)=\widehat K(\xi/r),
\]
and proves the corresponding exact fundamental-cell coefficient after
periodization.  It then constructs a volume-preserving regrouping of the six
variables into four transverse and two axial coordinates and proves
\[
 \int_{\mathbb R^6}\!\|\delta^5K(\delta x_\perp,x_{\rm ax})\|\,dx
   =\delta\int_{\mathbb R^6}\!\|K(x)\|\,dx,
 \qquad
 \widehat K_\delta(\xi)
   =\delta\widehat K(\xi_\perp/\delta,\xi_{\rm ax}).
\]
At aperture $\delta=1/N$ and shell scale $r=N^2$, the reciprocal argument is
machine-checked to be precisely the source/receiver chart
$(k_1/N,k_2/N,k_3/N^2)$ (and similarly for the second input).  Consequently
the periodized square-dyadic kernel has both the exact coefficient
\[
 \widehat K_N^{\rm per}(k,l)=\frac1N
 M_{\mathrm{norm},1/N}
 \left(\frac{k_1}{N},\frac{k_2}{N},\frac{k_3}{N^2};
       \frac{l_1}{N},\frac{l_2}{N},\frac{l_3}{N^2}\right)
\]
and the uniform fundamental-cell estimate
$\int\|K_N^{\rm per}\|\leq C/N$ for every $N\geq2$.
Thus the frozen-axis anisotropic square-dyadic transfer is discharged.  The
full-annulus and variable-frame assembly is described below.

The module \texttt{PancakeSquareDyadicFiniteTransfer} descends this final
kernel through the canonical half-open representative of the unit six-torus.
It proves that torus Fourier coefficients and torus $L^1$ mass are exactly the
corresponding fundamental-cell quantities, and hence that arbitrary finite
Fourier fields are acted on by the exact double sum of square-dyadic coherent
interactions.  The endpoint estimate is cardinality-free:
\[
 \|B_N(f,g)\|_{L^\infty(\mathbb T^3)}
 \leq 9\frac{C}{N}\,
   \|f\|_{L^\infty(\mathbb T^3)}
   \|g\|_{L^\infty(\mathbb T^3)},\qquad N\geq2.
\]
Thus the finite-field frozen-axis square-dyadic transfer is also discharged.

The module \texttt{PancakeFrozenFrameKernelTransfer} performs the frozen-frame
step at the kernel level.  Because a general rotation does not preserve
$\mathbb Z^3$, it first rotates the Euclidean cover kernel and only then
periodizes on the fixed standard lattice.  Orthogonal change of variables
preserves the exact cover $L^1$ mass, while Fourier covariance evaluates the
symbol at the rotated reciprocal frequency.  The module constructs a genuine
Euclidean linear isometry from every algebraic
\texttt{OrientedFrameEquiv}, lifts it diagonally to the two interacting
frequencies, and proves the exact finite-Fourier identity with the retained
coherent Biot--Savart interaction in the resulting oriented square-dyadic
chart.  One constant works for all such frames:
\[
 \|B_{N,F}(f,g)\|_{L^\infty(\mathbb T^3)}
 \leq 9\frac{C}{N}\,
   \|f\|_{L^\infty(\mathbb T^3)}
   \|g\|_{L^\infty(\mathbb T^3)},
 \qquad N\geq2.
\]
Thus frozen oriented-frame transference of the exact square-dyadic kernel is
discharged.  Physical vector conjugation, pointwise frame variation, and the
replacement of the initial demonstration chart by the full annulus are
handled next.

The module \texttt{PancakePhysicalFrameKernelTransfer} supplies the missing
input/output conjugation.  It complexifies the algebraic real frame, proves
that complexification commutes with inversion, and bounds both the forward and
inverse maps by the dimension-only coordinate constant $3$.  Applying the
frequency-rotated kernel to the two frame-coordinate inputs and rotating the
output back therefore gives
\[
 \|B^{\rm phys}_{N,F}(f,g)\|_{L^\infty}
 \leq 243\frac{C}{N}\|f\|_{L^\infty}\|g\|_{L^\infty}.
\]
The exact coherent finite-Fourier formula is proved with these three frame
maps in place.  More importantly, the same estimate is proved for arbitrary
continuous bounded fields and then for an arbitrary pointwise frame field
$F=F(x)$, simply by instantiating the frame-uniform kernel estimate at the
observation point.  No derivative of $F$ occurs.  This removes the
variable-frame commutator as an analytic loss at this particular pointwise
$L^\infty$ endpoint.

The modules \texttt{PancakeAnnularChartCutoff},
\texttt{PancakeAnnularUniformKernel}, and
\texttt{PancakeAnnularMultiplierAssembly} remove the remaining
demonstration-chart restriction.  On a square-dyadic annulus the cone
inequality forces each normalized axial coordinate into
$[-2,-3/4]\cup[3/4,2]$, while the transverse coordinates lie in $[-2,2]$.
There are therefore exactly four source/receiver sign combinations.  Smooth
compact buffers are constructed for them, their coherent denominators are
proved globally nonzero, and their inner cutoffs satisfy the exact partition
identity
\[
 \chi_{++}+\chi_{+-}+\chi_{-+}+\chi_{--}=1
\]
throughout the normalized annulus.  The parametric Schwartz estimate is
proved for every buffered cutoff and yields one aperture-uniform inverse
Fourier $L^1$ bound for the assembled four-chart kernel.  Its Fourier symbol
is machine-checked to equal the \emph{uncut} coherent Biot--Savart stretching
interaction on the whole annulus.

Finally, \texttt{PancakeAnnularFrozenFrameTransfer} and
\texttt{PancakeAnnularPhysicalFrameTransfer} carry this assembled multiplier
through the rotated Euclidean cover, square-dyadic rescaling, periodization,
physical input/output conjugation, and arbitrary pointwise frame fields.  The
exact finite-Fourier identity and the cardinality-free estimate
\[
 \|B^{\rm ann}_{N,F(x)}(f,g)(x)\|
 \le 243\frac{C}{N}\|f\|_{L^\infty}\|g\|_{L^\infty}
\]
hold uniformly for $N\ge2$.  Thus finite sign-chart coverage and partition
assembly are discharged.  The module
\texttt{PancakeAnnularSectorProjector} then defines the concrete finite
Fourier projector obtained by filtering a given support with the normalized
pancake-annulus predicate.  It proves, coordinate by coordinate, that the
six-dimensional frozen square-dyadic chart of a selected pair is exactly the
pair of their independently selected three-dimensional normalized modes.
Consequently the full-annulus premise of the exact physical kernel formula is
automatic on these projected fields; it is no longer a free support
hypothesis.  The subsequent module
\texttt{PancakeSquareDyadicSectorProjector} narrows this enclosing selector to
the genuine microlocal sector: radial frequency in $[N^2,2N^2)$ and squared
transverse aperture at most $N^{-2}$ after rotation into the chosen frame.
From these physical inequalities it proves the transverse bounds and the
nontrivial axial lower bound $3/4$, hence inclusion in the four-chart annulus.
The exact coherent-action formula therefore applies directly to the concrete
finite shell-and-cone projections.  The modules
\texttt{PancakeAxialFrameConstruction} and
\texttt{PancakeStrainAxialFrame} then remove the caller-supplied pointwise
frame.  The first constructs, by two explicit reflections, an
orientation-preserving orthogonal map sending any unit direction to the axial
vector; the second feeds it the smallest-eigenvalue direction supplied by the
ordered spectral theorem for the symmetric strain.  Thus the pointwise
$C/N$ endpoint now applies directly to a frame constructed from the strain.
Measurable dependence of the spectral choice through eigenvalue collisions,
and identification of these concrete finite projectors with the actual
Littlewood--Paley coherent share, are still required.

The modules \texttt{PancakeComplexMisalignmentExpansion},
\texttt{PancakeComplexCoherentPDEIdentification}, and
\texttt{PancakePhysicalMisalignmentSplit} close the algebraic split
behind that remaining analytic task.  The former expands the actual complex
Biot--Savart stretching interaction into the coherent--coherent term and the
three terms containing at least one misalignment factor, first modewise and
then under arbitrary finite source/receiver Fourier sums.  Thus the
finite-field norm analogue of \texttt{split\_le} follows from an exact
identity and one triangle inequality.  The latter proves that
$N^{-1}$ times the normalized annular symbol is literally the complexified
PDE stretching operator applied to the two coherent chart projections.  On
the genuine rotated square-dyadic sector, the exact physical kernel formula
is therefore rewritten directly as a sum of coherent--coherent PDE
interactions; no analogy between two separately defined expressions remains.
The third module performs the sum-level substitution: for arbitrary finite
Fourier data it proves the exact physical identity
\[
  \mathcal S_{N,F}
   =\mathcal K^{\mathrm{coh}}_{N,F}
    +\mathcal R^{\mathrm{mis}}_{N,F},
\]
where the remainder is precisely the sum of the coherent--misaligned,
misaligned--coherent, and misaligned--misaligned interactions.  A triangle
inequality and the annular kernel theorem give
\[
  |\mathcal S_{N,F}(x)|
   \leq 243(C/N)MP+|\mathcal R^{\mathrm{mis}}_{N,F}(x)|,
\]
with no dependence on the two Fourier support cardinalities.  The same theorem
is specialized to the pointwise frame built from the most compressive strain
eigenvector, so this statement no longer accepts a caller-supplied frame.  The
follow-up module \texttt{PancakeFullStrainFrame} fixes the residual rotation
around the compressive axis by an explicit orientation-preserving planar
rotation.  It proves that the inverse complex image of $e_x$ is exactly the
most expanding strain eigenvector and repeats the exact split and $C/N$ bound
in this full eigenline-aligned frame.  Consequently the Fourier remainder's
coherent reference is now the same expanding line used by the invariant
alignment-energy dynamics below.

The module \texttt{PancakeMisalignmentEnergyBridge} makes that compatibility
quantitative.  If a complex Fourier amplitude $u$ is divergence free at
frequency $k$, its exact coherent-line remainder $u-P_k u$ satisfies
\[
  |u-P_k u|^2
  \leq |u_y|^2+|u_z|^2
\]
in the full strain chart, where the chart's $x$-axis is the expanding
eigenline.  The proof applies the real orthogonal-projection inequality to
the real and imaginary parts and proves that complexification of the
oriented frame preserves the divergence pairing.  Summing over the actual
integer modes selected by the physical square-dyadic projector introduces no
cardinality factor.  An exact finite Parseval theorem then gives
\[
  \sum_{k\in\Gamma_{N,F}} |u_k-P_k u_k|^2
  \leq \int_{\mathbb T^3}
     \bigl(|U_y(x)|^2+|U_z(x)|^2\bigr)\,dx,
\]
where $U$ is the finite reconstruction in one frozen full strain frame.
This closes the coefficient-to-$L^2$ bridge; it does not by itself control a
frame varying with $x$, nor does an $L^2$ estimate alone supply the
$L^\infty$-based BKM integrand.

The follow-up module \texttt{PancakeTransverseEnergyFreezing} supplies the
local varying-frame comparison without choosing eigenvector signs.  For
symmetric strains $A,B$ with $\|A-B\|\leq\varepsilon$, it first proves the
top-eigenvalue Weyl bound
\[
  |\lambda_0(A)-\lambda_0(B)|\leq\|A-B\|.
\]
It then applies the matched-eigenline projector estimate.  For every positive
threshold, either a frozen top gap of $B$ is smaller than
$\text{threshold}+\varepsilon$, or
\[
  \sum_{k\in\Gamma_{N,A}} |u_k-P_{A,k}u_k|^2
  \leq \int_{\mathbb T^3}|U_{\perp B}(x)|^2\,dx
    +4\frac{\varepsilon}{\text{threshold}}
       \int_{\mathbb T^3}|U(x)|^2\,dx .
\]
Thus the eigenvector-sign and eigenvalue-collision problems are no longer
hidden in a continuity premise.  The remaining task is to aggregate this
cellwise alternative with the same-strain small-gap coercivity and the
space--time Littlewood--Paley budget.

The module \texttt{PancakeTransverseEnergyCoercivity} closes the separated
cell algebra.  If $S$ is symmetric, $e_0$ is its top eigenvector, and both top
gaps dominate $g$, then for every vector $u$---including $u=0$---
\[
 g\bigl(|u|^2-|u\mathbin{\cdot}e_0|^2\bigr)
 \leq \lambda_0(S)|u|^2-u\mathbin{\cdot}Su.
\]
The proof is the three-dimensional spectral expansion and is deliberately
homogeneous, so it never divides by vorticity magnitude.  Applying it to the
real and imaginary parts, integrating the finite reconstruction, and
composing with the freezing theorem yields the checked cellwise dichotomy:
either a frozen top gap is below $\mathrm{threshold}+\varepsilon$, or the
gap-weighted exact Fourier remainder is bounded by the integrated complex
Rayleigh defect plus
$(\mathrm{threshold}+\varepsilon)4\varepsilon/\mathrm{threshold}$ times the
full sector energy.  The remaining issue is no longer the separated-cell
estimate itself, but its bounded-overlap space--time aggregation and endpoint
scaling.

The finite combinatorics of that aggregation is now checked in
\texttt{PancakeCellAggregation}.  It packages one frozen comparison datum,
filters an arbitrary finite family into collision and separated cells, and
sums the separated coercive estimates with no support-cardinality constant.
Nonnegativity of both the complex Rayleigh defect and full energy allows the
right-hand sums to be enlarged to the whole family.  Two generic analytic
lemmas prove the corresponding bounded-overlap transfer: for measurable hard
cells of multiplicity at most $M$,
\[
  \sum_i\int_{Q_i}f\leq M\int f,
\]
and for nonnegative smooth weights $\chi_i$ satisfying
$\sum_i\chi_i\leq M$, the same conclusion holds for varying local densities
dominated by one global density.  This does \emph{not} yet identify the
global Parseval energy in each frozen comparison with a hard- or
smooth-localized cell energy.  A spatial cutoff changes Fourier support, so
that step requires a genuine localization/commutator or almost-orthogonality
estimate; bounded overlap alone cannot supply it.

The first part of that localization junction is now exact rather than an
interface.  The module \texttt{PancakeFiniteFourierLocalization} proves that
for finite Fourier cutoff $\chi$ and vector field $U$,
\[
 \chi(x)U(x)=\sum_{p,k}e^{2\pi i(p+k)\cdot x}\,\widehat\chi(p)\widehat U(k).
\]
Thus the localized output support is precisely contained in the Minkowski
sum of the cutoff and field supports, and every output mode differs from an
original sector mode by an actual cutoff frequency.  It also proves directly
in physical space that $\|\chi U\|_{L^2}^2\leq C^2\|U\|_{L^2}^2$ whenever
$\|\chi\|_{L^\infty}\leq C$; there is no mode-count loss.  What remains at
this junction is geometric: show that the permitted low cutoff frequencies
thicken the square-dyadic pancake sector by only a controlled margin, and
bound the mismatch between that thickened selector and the original smooth
projector.

There is a second field-level issue which a fixed-axis multiplier does not
address.  The actual pancake frame is the eigensystem of the low-frequency
strain and therefore depends on $(x,t)$.  The module
\texttt{PancakeFrameCovariance} proves that the continuous coherent estimate is
rotation-covariant for every frozen oriented orthonormal frame, including
arbitrary projected amplitudes.  The module
\texttt{PancakeVariableFrameCommutator} then gives the exact finite-Fourier
identity
\[
 T_{m(x,\cdot)}f=T_{m(x_0,\cdot)}f+
 T_{m(x,\cdot)-m(x_0,\cdot)}f
\]
and bounds the second term by the frequencywise symbol oscillation times the
actual Fourier-amplitude mass.  The uniform frozen-axis periodic kernel bound
is now supplied by \texttt{PancakeBilinearPeriodization}; exact finite-field
summation and the full anisotropic square-dyadic scaling are supplied by the
two subsequent modules above; frozen oriented-frame kernel transference is
supplied by \texttt{PancakeFrozenFrameKernelTransfer}; the full four-sign
annulus and its physical variable-frame endpoint are supplied by the annular
modules above, and the pointwise frame is now constructed from the strain.
The remaining field-transfer pieces are measurable control of that spectral
selection through eigenvalue collisions, aggregation of the proved
frozen-cell/small-gap alternative over a physical space--time cover, and
identification of the concrete finite-sector split with the actual
Littlewood--Paley time functionals.
The older frequencywise commutator estimate
remains correct, but the new uniform physical-kernel endpoint shows that no
frame-oscillation estimate is needed merely to obtain the pointwise
$L^\infty$ bound.
They are not silently included in the finite-sector estimate, and
\texttt{coherent\_le} has not yet been instantiated for the time-dependent
Navier--Stokes data structure.

\emph{The refined pin (committed).}  The module
\texttt{MisalignmentRefinedPin} splits the pancake-sector functional
into a direction-coherent share, carrying the interface form of the
proved aperture-gain kinematics, and a misalignment share; the new pin
\texttt{MisalignmentStrainBudget} demands the time-integrable envelope
only for the misalignment share, and the sealed reduction
\texttt{scaleLocalPancakeStrainBudget\_of\_misalignmentBudget}
reproduces the original budget with gain
$\mathrm{coherentGain} + \varepsilon$.  The refinement is strict: the
open burden moved from the full pancake stretching sum to its
misalignment share, whose controlling quantity is the direction
variation.  The vorticity-direction equation exposes a dynamical handle on
this quantity, but does not by itself give the required a priori estimate.

\emph{Exact unproven boundary.}  (i) Completion of the field-level transfer:
the modewise complex split, amplitude estimate, lattice count, exact
bilinear aperture factor, fixed-$\delta$ smooth tensor kernel, bilinear
scaling/Young bound, exact finite torus realization, frozen-frame covariance,
and variable/frozen split are now proved.  Uniform normalized Euclidean
bilinear-kernel mass, exact standard-six-torus Fourier realization, and the
frequency-uniform periodic $O(\delta)$ bilinear endpoint bound are also
proved.  Exact finite-Fourier summation, inner-chart identification with the
coherent PDE interaction, mass-preserving isotropic dyadic rescaling, exact
four-transverse/two-axial scaling, square-dyadic chart identification, its
periodized coefficient formula, its uniform $C/N$ mass bound, and uniform
frozen oriented-frame transference, physical input/output conjugation, and the
pointwise endpoint for an arbitrary frame field are now proved as well.
Four-sign coverage of the full normalized annulus, exact partition assembly,
one uniform inverse-Fourier mass bound for the assembled multiplier, its exact
uncut coherent-symbol identity on the annulus, and its physical
variable-frame $C/N$ endpoint are also proved.  The finite-support annular
sector projector is now constructed, its six-dimensional chart is identified
with the two independently filtered mode charts, and the genuine rotated
radial-shell/angle-cone selector is proved to lie in that annulus.  The exact
physical coherent formula follows without an additional support premise.  The
pointwise strain-derived axial frame and its expanding-line orientation are
also constructed.  Finally, the full
physical finite-sector stretching sum is proved to be exactly that annular
coherent action plus the three-term misalignment remainder, with a uniform
$C/N$ bound on the first term, including in the strain-derived frame.  The
complex coefficient remainder is now bounded by transverse expanding-line
energy, first modewise, then after finite-sector summation, and finally by an
exact physical-space $L^2$ integral via finite Parseval, with no
support-cardinality loss.  The remaining steps are measurable spectral
selection through eigenvalue collisions; a sign-invariant line-projector
theorem now controls the frozen-to-varying-frame error by
$4\varepsilon/\mathrm{threshold}$ times the full $L^2$ energy, with the
ordinary frozen top-gap band returned as the alternative.  On its separated
branch, homogeneous spectral-gap coercivity now bounds the transverse energy
by the integrated top Rayleigh defect.  The separated-cell finite sum and
generic bounded-overlap integration are now checked as well.  The remaining
steps are the localization/commutator estimate connecting global finite
reconstructions to the spatial cells.  Exact finite cutoff convolution,
Minkowski-sum support leakage, and the sharp physical $L^2$ cutoff bound are
now proved; controlled sector thickening and the projector commutator remain.
Further steps are aggregation of the small-gap branch
over space and time, conversion of this $L^2$/defect estimate to the required
scale-local endpoint budget, and identification and assembly of this concrete
finite-sector identity into \texttt{coherent\_le} for the actual
Navier--Stokes functionals.
(ii) The misalignment budget itself: dynamical,
time-integrated control of the misalignment share; the checked fences
of Sections~\ref{sec:lean}, \ref{sec:aperture} and the static no-go
apply to it verbatim.

\section{Opening the dynamical stage: direction rate and the tilting gap}
\label{sec:direction-dynamics}

The first dynamical audit separates two notions that should not be
conflated: rotation of the vorticity direction and rotation of the strain
eigenframe.  Both are relevant to a pancake configuration, but they obey
different equations.

For a nonzero vorticity vector written $\omega=r\xi$, with $|\xi|=1$, and a
material forcing $F$, define the transverse projection
\[
  P_{\xi}^{\perp}F=F-(\xi\cdot F)\xi,
  \qquad \dot\xi=r^{-1}P_{\xi}^{\perp}F.
\]
The module \texttt{DirectionEvolutionTilting} verifies this formula from a
differentiable factorization $\omega=r\xi$, including tangency
$\xi\cdot\dot\xi=0$.  If the pointwise vorticity equation is written as
$F=rS\xi+d$, where $S$ is the strain and $d$ is the diffusive remainder,
the exact split is
\[
  \dot\xi=P_{\xi}^{\perp}(S\xi)
       +r^{-1}P_{\xi}^{\perp}d.
\]
For two differentiable direction paths, the same module proves
\[
  \frac{d}{dt}|\xi_1-\xi_2|^2
   =2(\xi_1-\xi_2)\cdot(\dot\xi_1-\dot\xi_2),
\]
together with its Cauchy--Schwarz bound by
$2|\xi_1-\xi_2|(|\dot\xi_1|+|\dot\xi_2|)$.  Thus total direction turning is
an exact dynamical input to misalignment production; obtaining the
scale-critical time integral required by \texttt{MisalignmentStrainBudget}
remains analytic work.

The pressure Hessian enters through the strain-eigenframe equation instead.
Differentiating $Se_i=\lambda_i e_i$ and pairing with $e_j$, $i\ne j$, gives
\[
  \langle e_j,\dot e_i\rangle
   =\frac{\langle e_j,\dot S e_i\rangle}{\lambda_i-\lambda_j}.
\]
The anisotropic pressure Hessian contributes to the numerator.  The checked
theorem \texttt{no\_uniformGapFreePressureTiltingBound} records the necessary
denominator: no nonnegative constant bounds the angular rate from a bounded
off-diagonal torque independently of the strain eigenvalue gap.  Unit torque
with gap $1/(C+1)$ already defeats a proposed constant $C$.

This rules out a raw pressure-Hessian norm as the next envelope.  A viable
dynamical continuation must instead control a gap-weighted pressure torque,
or prove cancellation/residence-time control over near-degenerate intervals,
and then transfer that control to the Littlewood--Paley misalignment share.

The same denominator is now proved for spatial frame freezing.  In
\texttt{PancakeEigenframePerturbation}, if $Ae=\lambda e$, $Bf=\mu f$, $B$
is self-adjoint, and $e,f$ are unit vectors, then
\[
  (\lambda-\mu)\langle f,e\rangle
    =\langle f,(A-B)e\rangle,
  \qquad
  |\lambda-\mu|\,|\langle f,e\rangle|\leq\|A-B\|.
\]
Consequently, at every positive threshold, either the eigenvalues are in the
near-degenerate band or the frame overlap is bounded by the strain oscillation
divided by that threshold.  This is the exact gap/near-degeneracy alternative
needed by the variable-frame remainder; it is not yet the scale-summed bound
needed for closure.

Two further modules turn this into an intrinsic local dynamics statement.
\texttt{PancakeStrainSpectralFrame} constructs, for every symmetric
three-dimensional strain operator, an ordered orthonormal eigenframe and
proves the three-coordinate Parseval and Rayleigh identities.  Thus the
spectral coordinates used below are constructed objects, not assumed data.
For a unit vorticity direction $\xi$ and a moving unit top eigenvector $e$,
\texttt{PancakeAlignmentDynamics} defines
\[
  E(\xi,e)=1-\langle\xi,e\rangle^2
\]
and proves the exact simple-gap identity
\[
 \dot E
 =-2\langle\xi,e\rangle^2
     \bigl(\lambda_0-\langle\xi,S\xi\rangle\bigr)
   -2\langle\xi,e\rangle
     \bigl(\langle d,e\rangle+\langle\xi,\theta\rangle\bigr),
\]
where $d$ is the non-strain direction forcing and $\theta=\dot e$.  The
Rayleigh defect is bounded below by the top spectral gap times $E$, and the
module integrates this damping to obtain the residence estimate
\[
 \int_0^T (\lambda_0-\lambda_1)E\,dt
 \le E(0)+2\int_0^T
   |\langle d,e\rangle+\langle\xi,\theta\rangle|\,dt
\]
under its stated regularity and simple-gap hypotheses.  It also proves the
complementary upper bound by the full spectral width.

These identities expose rather than hide the hard cases.  Substituting the
usual eigenvector perturbation estimate
$|\theta|\lesssim|\dot S|/(\lambda_0-\lambda_1)$ into the residence inequality
and applying a naive Young inequality produces still worse inverse-gap costs;
it is not a closure.  Moreover, a large trace-free strain may be isotropic on
a strongly expanding two-plane.  Such a partially degenerate plane supplies
damping toward the plane but does not select a coherent line inside it.
Controlling those intervals requires a projector/cluster formulation or a
genuine time-integrated cancellation, and remains part of the dynamical pin.

\section{Status and scope}
\label{sec:scope}

We state the scope of this finding exactly.

\paragraph{What is established from the original audit.} Hypothesis~(H1) of
v7 --- a uniform
$\Hm\to\Hm$ operator bound \eqref{eq:H1} on the adjoint / push-forward $\Ad_g$
over an $\Hm$ chart-neighborhood of the identity --- is \emph{false} in the
normalized Fourier-mode shear model of Section~\ref{sec:lean}. The refutation is
machine-checked (no \texttt{sorry}, elementary) in
\texttt{Mettapedia.FluidDynamics.NavierStokes.BenH1Break}, and the underlying
norms are independently derived from the Sobolev-weight structure in the
companion numeric lab. The model faithfully captures the mechanism that defeats
v7's local-boundedness argument: the adjoint \eqref{eq:adg} loses one derivative
($g\in\Hm\Rightarrow Dg\in H^{m-1}$), so on a fixed-radius $\Hm$ chart the
operator norm is unbounded along high-wavenumber shears.

\paragraph{What is not established.} This is a refutation of (H1)
\emph{in a model} --- a single analytic Fourier-mode family together with the
linearized commutator action --- not a from-Mathlib-Sobolev-first-principles
theorem about $\SDiff^m(\mathbb{T}^3)$. The model uses the exact push-forward
formula \eqref{eq:adg} on an explicit shear family and the standard Fourier
$\Hm$ norm; it does not formalize the full diffeomorphism-group functional
analysis. The mechanism (derivative loss) is captured faithfully, but the
strength of the conclusion is ``checked in this model,'' not ``checked over the
entire chart from Mathlib's Sobolev theory.''

\paragraph{Consequence for v7.} Because every quantitative step of the
SG--Cole--Hopf program inherits the constant $\CAd$ from \eqref{eq:H1} ---
through \eqref{eq:prop65}, \eqref{eq:logbound}, and \eqref{eq:vort} --- the
falsity of~(H1) as stated means the route does \emph{not}, as written, establish
global regularity. v7 reduces 3D Navier--Stokes global regularity to~(H1)
(equivalently, to a uniform curvature / Bakry--\'Emery-type bound), and that
reduction target does not hold uniformly in the form required.

\paragraph{What the repair establishes.} The stochastic-Lagrangian repair is
not refuted by the H1 countermodel, and that evasion is itself a checked
theorem (Section~\ref{sec:evasion}).  Machine-checked without \texttt{sorry}
across the campaign: the one-form calculus with the exact frame-Laplacian
identity; the divergence-free non-exact witness forcing the gauge; the
pushdown theorem producing Navier--Stokes with the explicit pressure
$p=\partial_tW+\tfrac12\langle u,u\rangle+\langle u,\nabla W\rangle-\nu\Delta W$
from the averaged transport hypothesis; the gate consumers and the assembled
conditional route; nonzero exact inhabitants (decaying shear with balanced
gauge) and end-to-end rest canaries for both the coarse and the refined
route; the frozen-strain exact solution with its unbounded-amplification and
coverage-gap theorems; the aperture dichotomy--budget equivalence with both
gap-band failure theorems and the Gr\"onwall-saturation bound; the
single-mode Biot--Savart null structure; and the dyadic absorption algebra
with explicit constants.  The theorem
\texttt{stochasticLagrangian\_conditional\_route\_of\_pancakeBudget} records
that the route closes if the scale-local pancake strain budget and the
non-pancake sector estimates are supplied.

\paragraph{Scope of the new pieces.}  The same honesty discipline as for the
H1 model applies.  The frozen-strain results are theorems about a linear
prescribed-strain model, labeled as such; the aperture theorems are exact
inequalities between the model's threshold parameters; the plane-wave null
structure is an identity for one exact Fourier mode (sector aggregation is
exactly the open budget); the averaged transport equation enters the
pushdown as a hypothesis field (no stochastic calculus is formalized); and
the coercive/residual sector estimates are named interfaces consumed from
the source analysis, not re-proved.  None of these statements is a claim
about general Navier--Stokes solutions beyond what its label says.

\paragraph{Verification metadata.}  The named campaign modules contain no
\texttt{sorry}, \texttt{admit}, or unsafe definitions.  Selected load-bearing
theorems carry explicit machine axiom audits; the remaining newly added
headlines, including the finite Parseval misalignment-energy bridge, are
checked by the Navier--Stokes live surface and the same explicit audit
boundary.  The live surface builds green.  Earlier
conditional route theorems retain their instantiation canaries demonstrating
that their stated hypotheses are simultaneously satisfiable.  This is a
claim about proof integrity, not about discharge of the explicitly open
analytic hypotheses.

\paragraph{What remains open.} We make no claim to have proved the Millennium
Navier--Stokes statement.  The latest pancake-closure proposal fails as a
complete proof because its aperture split requires the pointwise scale-local
condition $\gamma\le\nu\lambda_q$, and that condition is exactly the burden of
the unresolved self-consistent strain estimate.  After the September 2026
refinement (Section~\ref{sec:coherence}), the open analytic input is
\texttt{MisalignmentStrainBudget} together with the field-level
coherent/misaligned transfer, which jointly reproduce
\texttt{ScaleLocalPancakeStrainBudget} (cross-referenced with the coarser
\texttt{StochasticStretchingEstimate}).  A successful completion must prove
measurable control of the pointwise strain-derived spectral frame through
eigenvalue collisions, identify the now-exact finite-sector physical split
with the actual scale-local Littlewood--Paley time functionals, aggregate the
proved frozen-cell/small-gap projector alternative and convert its $L^2$
control to the endpoint scale budget, and
prove a dynamical,
time-integrated control of the vorticity-direction misalignment share, strong
 enough to cover the gap band identified above.  The finite four-sign
annular partition and its exact periodic physical transference are proved;
the remaining identification may not assign the aperture gain to the
projection.
Section
\ref{sec:direction-dynamics} further shows that a pressure-Hessian route must
carry inverse strain-eigenvalue-gap information or an equivalent
time-integrated cancellation; pressure-Hessian size alone cannot close the
pin.  Once the analytic estimates are closed, the remaining endpoint work is
to promote the common pointwise majorant to the measurable spatial
essential-supremum BKM integrand and apply a formal continuation theorem for
the chosen solution class.

\section*{Provenance}
The manuscript discussed is Ben Goertzel, \emph{Global Regularity of 3D
Navier--Stokes via SG--Cole--Hopf Program}, version v7, together with the
subsequent pancake-closure repair proposal audited in July 2026.  The
formalization, Fourier-mode model, repaired conditional route, pancake
adjudication, and numeric lab are work of the MeTTapedia formalization
effort.  The repaired-route development comprises, under
\texttt{Mettapedia.FluidDynamics.NavierStokes}: \texttt{BenH1Break},
\texttt{OneFormFrameCalculus}, \texttt{OneFormExactness},
\texttt{GradientOneForm}, \texttt{FrozenStrainModel},
\texttt{ApertureCoverageDichotomy}, \texttt{PlaneWaveSelfStrain}, and the
\texttt{StochasticLagrangian} layer (\texttt{PullbackOneForm},
\texttt{ShearJacobianBounds}, \texttt{ConstantinIyerRepresentation},
\texttt{VorticityStretchingGate}, \texttt{DecayingShearInstance},
\texttt{DyadicPancakeClosure}, \texttt{PancakeCoherentPairEstimate},
\texttt{PancakeRealCoherentPairEstimate}, \texttt{PancakeConeOperatorKernel},
\texttt{PancakePeriodicCoherentSplit},
\texttt{PancakePeriodicComplexStretch},
\texttt{PancakePeriodicOperatorKernel}, \texttt{PancakeFrameCovariance},
\texttt{PancakeVariableFrameCommutator}, and
\texttt{PancakeEigenframePerturbation},
\texttt{PancakeAlignmentDynamics}, \texttt{PancakeStrainSpectralFrame},
\texttt{PancakeSmoothCoherentSymbol},
\texttt{PancakeSmoothCoherentAperture},
\texttt{PancakeCoherentStretchFactorization},
\texttt{PancakeSmoothCoherentStretchSymbol}, and
\texttt{PancakeBilinearKernelScaling}, and
\texttt{PancakeUniformKernelFamily}, and
\texttt{PancakeBilinearPeriodization},
\texttt{PancakeFiniteFourierBilinearTransfer},
\texttt{PancakeDyadicKernelScaling},
\texttt{PancakeSquareDyadicFiniteTransfer}, and
\texttt{PancakeFrozenFrameKernelTransfer}, and
\texttt{PancakePhysicalFrameKernelTransfer}), each with a regression companion where
applicable.  The aperture module is retained as an explicitly attenuated
analytic model; the exact coherent projection no-go above prevents its use as
the desired field-transfer estimate.
Background references: J.~T. Beale, T.~Kato, A.~Majda, \emph{Comm.\ Math.\
Phys.}\ \textbf{94} (1984), 61--66; P.~Constantin, G.~Iyer, \emph{Comm.\
Pure Appl.\ Math.}\ \textbf{61} (2008), 330--345; D.~Ebin, J.~Marsden,
\emph{Ann.\ of Math.}\ \textbf{92} (1970), 102--163; H.~Bahouri,
J.-Y.~Chemin, R.~Danchin, \emph{Fourier Analysis and Nonlinear Partial
Differential Equations}, Springer, 2011; P.~Constantin, C.~Fefferman,
\emph{Indiana Univ.\ Math.\ J.}\ \textbf{42} (1993), 775--789; J.~Tom,
M.~Carbone, A.~D.~Bragg, \emph{J.\ Fluid Mech.}\ \textbf{910} (2021), A24.

\end{document}
